%! Parametric Functions Modeling Planar Motion — Unit 4, Lesson 4.2 — Homework (Answer Key)
\documentclass[10pt]{article}
\usepackage{precalculus-article}
\usepackage{precalculus-key}   % pulls in -boxes; do NOT also load -boxes
\newcommand{\TallMath}[1]{$\displaystyle #1\rule[-1.4em]{0pt}{3.2em}$}

\begin{document}

\pageheader{Unit 4, Lesson 4.2}{Homework --- Answer Key}
\namedateperiod

\vspace{0.10in}

\begin{notesbox}{Parametric Planar Motion Practice --- Key}

\textbf{1.}\quad Let $f(t) = (t^2 - 2t,\; 3t - t^2)$ for $0 \le t \le 4$.

\vspace{4pt}
(a)\quad Table:

\vspace{4pt}
\begin{center}
\renewcommand{\arraystretch}{1.6}
\begin{tabular}{|c|c|c|c|}
\hline
$t$ & $x(t) = t^2 - 2t$ & $y(t) = 3t - t^2$ & $(x,\; y)$ \\
\hline
$0$ & \ans{$0$}  & \ans{$0$}  & \ans{$(0,\;0)$}  \\
\hline
$1$ & \ans{$-1$} & \ans{$2$}  & \ans{$(-1,\;2)$} \\
\hline
$2$ & \ans{$0$}  & \ans{$2$}  & \ans{$(0,\;2)$}  \\
\hline
$3$ & \ans{$3$}  & \ans{$0$}  & \ans{$(3,\;0)$}  \\
\hline
$4$ & \ans{$8$}  & \ans{$-4$} & \ans{$(8,\;-4)$} \\
\hline
\end{tabular}
\end{center}

\vspace{8pt}
(b)\quad Horizontal extrema:

\ans{$x(t) = t^2 - 2t$; vertex at $t = 1$. $x(1) = -1$ (minimum,
leftmost position).\\
Endpoints: $x(0) = 0$, $x(4) = 8$ (rightmost). Maximum is $x = 8$.}

\vspace{0.08in}
(c)\quad Vertical extrema:

\ans{$y(t) = 3t - t^2 = -(t^2 - 3t)$; vertex at $t = 3/2$.\\
$y(3/2) = 9/2 - 9/4 = 9/4 = 2.25$ (maximum, highest position).\\
Endpoints: $y(0) = 0$, $y(4) = -4$ (minimum, lowest position).}

\vspace{0.10in}
\textbf{2.}

\vspace{4pt}
(a)\quad $x$-intercepts (zeros of $y(t)$):

\ans{$3t - t^2 = t(3-t) = 0 \;\Rightarrow\; t = 0$ or $t = 3$.\\
$x(0) = 0$, $x(3) = 3$. $x$-intercepts: $(0,\,0)$ and $(3,\,0)$.}

\vspace{0.04in}
(b)\quad $y$-intercepts (zeros of $x(t)$):

\ans{$t^2 - 2t = t(t-2) = 0 \;\Rightarrow\; t = 0$ or $t = 2$.\\
$y(0) = 0$, $y(2) = 2$. $y$-intercepts: $(0,\,0)$ and $(0,\,2)$.}

\end{notesbox}

\vspace{0.08in}

\begin{notesbox}{Contextual Problem --- Key}

\textbf{3.}\quad $x(t) = 4t$, $y(t) = -5t^2 + 20t$, $0 \le t \le 4$.

\vspace{4pt}
(a)\quad Maximum height:

\ans{$y(t) = -5t^2 + 20t$; vertex at $t = -20/(2 \cdot (-5)) = 2$ s.\\
$y(2) = -20 + 40 = 20$ m. The ball reaches a maximum height of
$20$ meters at $t = 2$ seconds.}

\vspace{0.04in}
(b)\quad When does the ball hit the ground?

\ans{$-5t^2 + 20t = 0 \;\Rightarrow\; -5t(t - 4) = 0 \;\Rightarrow\;
t = 4$ s (excluding $t = 0$).\\
$x(4) = 4 \cdot 4 = 16$ m east. The ball lands at $x = 16$ m east.}

\end{notesbox}

\vspace{0.08in}

\begin{notesbox}{AP-Style Practice --- Key}

\textbf{4.}\quad $h(t) = (t^2 - 2t - 3,\; t^2 - 4)$,
$-1 \le t \le 4$. Which correctly identifies the $y$-intercepts?

\vspace{6pt}
\begin{enumerate}[label=(\Alph*), leftmargin=2em, itemsep=4pt]
  \item $(-4,\; 0)$ and $(5,\; 0)$
  \item \textcolor{keyred}{\textbf{$(0,\; -3)$ and $(0,\; 5)$ $\leftarrow$ correct}}
  \item $(0,\; -3)$ only
  \item $(0,\; -4)$ and $(0,\; 12)$
\end{enumerate}

\begin{teachernote}
  \textbf{Solution:} $y$-intercepts occur where $x(t) = 0$.\\
  $t^2 - 2t - 3 = (t-3)(t+1) = 0 \;\Rightarrow\; t = 3$ or $t = -1$.
  Both are in $[-1, 4]$.\\
  At $t = -1$: $y(-1) = 1 - 4 = -3$ → $(0,\,-3)$.\\
  At $t = 3$: $y(3) = 9 - 4 = 5$ → $(0,\,5)$.\\
  Answer: \textbf{(B)}.

  \textbf{Distractor analysis:}\\
  (A) $(-4, 0)$ and $(5, 0)$: student found $x$-intercepts (zeros of
  $y(t)$) instead of $y$-intercepts --- the classic rule-reversal error.\\
  (C) $(0, -3)$ only: student excluded $t = -1$ because they treated
  the domain as open or forgot to check all zeros.\\
  (D) Evaluated $y$ at wrong $t$-values or made arithmetic errors.

  Post the rule on the board: \emph{zeros of $x(t)$ → $y$-intercepts
  of the curve; zeros of $y(t)$ → $x$-intercepts.} If most students
  choose (A), spend two minutes restating this rule before moving on.
\end{teachernote}

\end{notesbox}

\end{document}
